%! Dividing Polynomials; Remainder and Factor Theorems — Unit 4, Lesson 4.3 — Group Activity (Answer Key)
\documentclass[10pt]{article}
\usepackage{algebra2-article}
\usepackage{algebra2-key}   % pulls in -boxes; do NOT also load -boxes
\newcommand{\TallMath}[1]{$\displaystyle #1\rule[-1.4em]{0pt}{3.2em}$}
\newcommand{\mans}[1]{{\color{keyred}\boldsymbol{#1}}}   % math-mode answer (\ans is text-only)

\begin{document}

\pageheader{Unit 4, Lesson 4.3}{Group Activity --- Answer Key}
\namepartnerperiod

\vspace{0.08in}

\begin{headlinebox}{forestbg}
\small\textbf{Divide and Conquer.} Every problem today runs on one statement:
\textbf{dividend $=$ divisor $\times$ quotient $+$ remainder}. Three habits keep you out of trouble:
write the dividend in standard form with a \emph{placeholder $0$} for every missing power; remember
that $x + 2$ means $r = -2$; and read the remainder --- a remainder of $0$ is the Factor Theorem
telling you that you just found a factor.
\end{headlinebox}

\vspace{0.06in}

\begin{tcolorbox}[colback=white, colframe=black!40, title=\textbf{Tier R --- Remediate},
    fonttitle=\bfseries, arc=1.5mm, left=3mm, right=3mm, top=2mm, bottom=2mm, breakable]
\begin{enumerate}[leftmargin=1.7em, itemsep=9pt]
  \item \textbf{Monomial divisor.} Split the fraction term by term.

    \vspace{2pt}
    $\dfrac{10x^4 - 15x^3 + 5x^2}{5x^2} = $ \ans{$2x^2 - 3x + 1$}

    \vspace{2pt}
    Careful with the last term --- it divides out to \ans{$1$}, not $0$.

  \item \textbf{Long division, step by step.} Divide $x^2 + 7x + 12$ by $x + 3$.

    \vspace{3pt}
    \emph{Step 1 --- first term of the quotient:} $\dfrac{x^2}{x} = $ \ans{$x$}

    \vspace{3pt}
    \emph{Step 2 --- multiply and subtract:} $x(x+3) = $ \ans{$x^2+3x$}, and
    $(x^2+7x) - ($\ans{$x^2+3x$}$) = $ \ans{$4x$}

    \vspace{3pt}
    \emph{Step 3 --- bring down the $12$ and repeat:} $\dfrac{4x}{x} = $ \ans{$4$}, \;
    $4(x+3) = $ \ans{$4x+12$}

    \vspace{3pt}
    \emph{Step 4:} Quotient \ans{$x+4$}, remainder \ans{$0$}. \; Check by multiplying:
    $(x+3)($\ans{$x+4$}$) = x^2+7x+12$? \ans{Yes}

  \item \textbf{Synthetic division.} Divide $x^3 + 3x^2 - 4x - 12$ by $x + 3$.

    \vspace{3pt}
    First, the sign flip: the divisor is $x + 3$, so $r = $ \ans{$-3$}
    \begin{center}
    \renewcommand{\arraystretch}{1.6}
    \begin{tabular}{r|C{1.4cm}C{1.4cm}C{1.4cm}C{1.4cm}}
    \ans{$-3$} & $1$ & $3$ & $-4$ & $-12$ \\
        &     & \ans{$-3$} & \ans{$0$} & \ans{$12$} \\ \cline{2-5}
        & \ans{$1$} & \ans{$0$} & \ans{$-4$} & \ans{$0$} \\
    \end{tabular}
    \end{center}

    \vspace{2pt}
    Quotient: \ans{$x^2-4$} \quad Remainder: \ans{$0$}

    \vspace{2pt}
    The quotient is a difference of squares, so factor completely:
    $x^3+3x^2-4x-12 = $ \ans{$(x+3)(x-2)(x+2)$}

  \item \textbf{Remainder Theorem check.} For that same $f(x) = x^3 + 3x^2 - 4x - 12$, compute
    $f(-3) = $ \ans{$0$}

    \vspace{2pt}
    Compare it with the remainder in problem 3. What do you notice? \ans{They are the same --- that is the Remainder Theorem.}
\end{enumerate}
\end{tcolorbox}

\begin{tcolorbox}[colback=white, colframe=black!40, title=\textbf{Tier A --- Approaching Proficiency},
    fonttitle=\bfseries, arc=1.5mm, left=3mm, right=3mm, top=2mm, bottom=2mm, breakable]
\begin{enumerate}[leftmargin=1.7em, itemsep=10pt]
  \item \textbf{Long division with a gap.} Divide $x^3 - 5x + 3$ by $x - 2$. Write the dividend with
    its placeholder first.

    \vspace{2pt}
    Dividend in standard form: \ans{$x^3 + 0x^2 - 5x + 3$}

    \[
    \begin{array}{r}
      \mans{x^2+2x-1}\phantom{{}+3} \\[2pt]
      x-2\,\overline{\smash{\big)}\,x^3+0x^2-5x+3} \\[2pt]
      \underline{-(x^3-2x^2)}\phantom{{}-5x+3} \\[2pt]
      2x^2-5x\phantom{{}+3} \\[2pt]
      \underline{-(2x^2-4x)}\phantom{{}+3} \\[2pt]
      -x+3 \\[2pt]
      \underline{-(-x+2)} \\[2pt]
      1
    \end{array}
    \]

    Quotient: \ans{$x^2+2x-1$} \quad Remainder: \ans{$1$} \quad
    $\dfrac{x^3-5x+3}{x-2} = $ \ans{$x^2+2x-1+\dfrac{1}{x-2}$}

  \item \textbf{Synthetic division, then finish the factoring.} Divide $2x^3 + 5x^2 - x - 6$ by
    $x + 2$.
    \begin{center}
    \renewcommand{\arraystretch}{1.6}
    \begin{tabular}{r|C{1.4cm}C{1.4cm}C{1.4cm}C{1.4cm}}
    \ans{$-2$} & \ans{$2$} & \ans{$5$} & \ans{$-1$} & \ans{$-6$} \\
        &     & \ans{$-4$} & \ans{$-2$} & \ans{$6$} \\ \cline{2-5}
        & \ans{$2$} & \ans{$1$} & \ans{$-3$} & \ans{$0$} \\
    \end{tabular}
    \end{center}

    \vspace{2pt}
    Depressed polynomial: \ans{$2x^2+x-3$} \quad Remainder: \ans{$0$}

    \vspace{2pt}
    Factor the quotient, then write the complete factorization: \ans{$(x+2)(2x+3)(x-1)$}

    \vspace{2pt}
    The three zeros are \ans{$-2,\ -\frac{3}{2},\ 1$} \; --- one of them is \emph{not} an integer. Which?
    \ans{$-\frac{3}{2}$}

  \item \textbf{No dividing allowed.} Let $g(x) = x^4 - 3x^2 + 2x - 8$. Use the Remainder Theorem
    only.

    \vspace{2pt}
    Remainder when divided by $(x-2)$: \ans{$0$} \qquad by $(x+1)$: \ans{$-12$}

    \vspace{2pt}
    Which of the two is a factor of $g$, and how do you know? \ans{$(x-2)$ --- its remainder $g(2)$ is $0$, which is the Factor Theorem.}

  \item \textbf{Factor Theorem sweep.} Let $p(x) = x^3 - x^2 - 9x + 9$. Test each candidate.

    \vspace{2pt}
    $p(1) = $ \ans{$0$} \quad $p(-1) = $ \ans{$16$} \quad $p(3) = $ \ans{$0$} \quad
    $p(-3) = $ \ans{$0$}

    \vspace{2pt}
    How many factors of the form $(x-r)$ did you just find? \ans{$3$} \;
    Write $p$ completely factored: \ans{$(x-1)(x-3)(x+3)$}

    \vspace{2pt}
    This one also factors by \emph{grouping} (Lesson 4.2). Show that route in one line:
    \ans{$x^2(x-1)-9(x-1)=(x-1)(x^2-9)$}

  \item \textbf{Error analysis.} Two setups, two different mistakes.

    \vspace{2pt}
    (a) To divide $x^3 - 4x + 1$ by $x - 2$, a student writes the top row as \; $2 \mid 1 \;\; -4
    \;\; 1$. \; What is missing? \ans{the placeholder $0$ for $x^2$}

    \vspace{2pt}
    Write the correct top row: \ans{$2 \mid 1\;\; 0\;\; -4\;\; 1$} \; and finish: quotient \ans{$x^2+2x$}, remainder
    \ans{$1$}

    \vspace{2pt}
    (b) To divide by $x + 5$, a student puts $5$ in the corner box. Why is that wrong, and what
    belongs there? \ans{Synthetic division uses $r$ from $x-r$; since $x+5=x-(-5)$, the corner takes $-5$.}
\end{enumerate}
\end{tcolorbox}

\begin{tcolorbox}[colback=white, colframe=black!40, title=\textbf{Tier E --- Extension},
    fonttitle=\bfseries, arc=1.5mm, left=3mm, right=3mm, top=2mm, bottom=2mm, breakable]
\textbf{Part 1 --- A divisor synthetic division cannot touch.} Divide
$x^4 + 2x^3 - 7x^2 - 8x + 12$ by $x^2 + x - 6$.
\begin{enumerate}[label=(\alph*), leftmargin=1.8em, itemsep=6pt]
  \item Why can this one \emph{not} be done synthetically? \ans{Synthetic division only works for divisors of the form $x-r$; this divisor is quadratic.}
  \item Long-divide. \quad Quotient: \ans{$x^2+x-2$} \quad Remainder: \ans{$0$}

    \[
    \begin{array}{r}
      \mans{x^2+x-2}\phantom{{}-8x+12} \\[2pt]
      x^2+x-6\,\overline{\smash{\big)}\,x^4+2x^3-7x^2-8x+12} \\[2pt]
      \underline{-(x^4+\phantom{1}x^3-6x^2)}\phantom{{}-8x+12} \\[2pt]
      x^3-\phantom{1}x^2-8x\phantom{{}+12} \\[2pt]
      \underline{-(x^3+\phantom{1}x^2-6x)}\phantom{{}+12} \\[2pt]
      -2x^2-2x+12 \\[2pt]
      \underline{-(-2x^2-2x+12)} \\[2pt]
      0
    \end{array}
    \]

  \item Both the divisor and the quotient factor. Write the original polynomial as a product of four
        linear factors: \ans{$(x+3)(x-2)(x+2)(x-1)$}
  \item List its four zeros: \ans{$-3,\ 2,\ -2,\ 1$}
\end{enumerate}

\vspace{5pt}
\textbf{Part 2 --- Solve for the missing coefficient.} Find the value of $k$ that makes $(x-3)$ a
factor of $f(x) = x^3 + kx^2 - 5x + 6$.
\begin{enumerate}[label=(\alph*), leftmargin=1.8em, itemsep=6pt]
  \item By the Factor Theorem, $(x-3)$ is a factor exactly when $f(3) = $ \ans{$0$}. Write that
        equation in terms of $k$: \ans{$27+9k-15+6=0$, i.e.\ $9k+18=0$}
  \item Solve: $k = $ \ans{$-2$} \; So $f(x) = $ \ans{$x^3-2x^2-5x+6$}
  \item Divide out $(x-3)$ and factor completely: \ans{$(x-3)(x+2)(x-1)$}
\end{enumerate}

\vspace{5pt}
\textbf{Part 3 --- Prove it, then exploit it.}
\begin{enumerate}[label=(\alph*), leftmargin=1.8em, itemsep=6pt]
  \item Start from $f(x) = (x-r)q(x) + R$ and substitute $x = r$ to prove the Remainder Theorem in
        two lines. \ansline{$f(r) = (r-r)q(r) + R = 0\cdot q(r) + R = R$. So the remainder $R$ equals $f(r)$ --- always.}
  \item Because the remainder \emph{is} $f(r)$, synthetic division is also a way to \emph{evaluate} a
        polynomial. Use it to find $f(4)$ for $f(x) = x^4 - 3x^3 + 2x - 7$:
    \begin{center}
    \renewcommand{\arraystretch}{1.6}
    \begin{tabular}{r|C{1.3cm}C{1.3cm}C{1.3cm}C{1.3cm}C{1.3cm}}
    \ans{$4$} & \ans{$1$} & \ans{$-3$} & \ans{$0$} & \ans{$2$} & \ans{$-7$} \\
        &     & \ans{$4$} & \ans{$4$} & \ans{$16$} & \ans{$72$} \\ \cline{2-6}
        & \ans{$1$} & \ans{$1$} & \ans{$4$} & \ans{$18$} & \ans{$65$} \\
    \end{tabular}
    \end{center}
    $f(4) = $ \ans{$65$} \; Now check by substituting $4$ directly: \ans{$256-192+8-7=65$}
  \item Which method used fewer multiplications? \ans{synthetic division (4, versus building each power)}
\end{enumerate}

\vspace{5pt}
\textbf{Part 4 --- Tomorrow's problem, today.} Nobody will hand you the value of $r$ forever. Take
$h(x) = 2x^3 - 3x^2 - 11x + 6$ and test the integers that divide the constant term $6$.
\begin{enumerate}[label=(\alph*), leftmargin=1.8em, itemsep=6pt]
  \item Candidates: $\pm 1, \pm 2, \pm 3, \pm 6$. Which two give $h(r) = 0$? \ans{$r=-2$ and $r=3$}
  \item Divide out one of them, say $(x+2)$, and factor the quotient completely:
        $h(x) = $ \ans{$(x+2)(2x-1)(x-3)$}
  \item List all three zeros: \ans{$-2,\ \frac{1}{2},\ 3$} \; One of them was \emph{not} on your candidate list.
        Which, and why did the list miss it? \ansline{$\frac{1}{2}$. The list only held integers dividing the constant term $6$, and this zero is a fraction --- its denominator $2$ comes from the leading coefficient.}
  \item So what would a complete candidate list have to include? \ansline{Fractions: every $\frac{p}{q}$ with $p$ a factor of the constant term and $q$ a factor of the leading coefficient. That is the Rational Root Theorem --- Lesson 4.4.}
\end{enumerate}
\end{tcolorbox}

\begin{teachernote}
\textbf{Tier R} is procedure with the theorem attached: item~4 exists only so students see the
Remainder Theorem confirm the number they just computed in item~3. Do not let them skip it.
\textbf{Tier A} items 1 and 5(a) are the same error from two sides --- the missing placeholder ---
and 5(b) is the sign flip; between them they account for most of the wrong answers you will collect
this week. Item~2's non-integer zero $-\frac{3}{2}$ is deliberate foreshadowing; item~4's second
route rewards students who remember grouping from Lesson 4.2.

\vspace{3pt}
\textbf{Tier E, Part 1} is the one item in the packet that satisfies the ``factorable trinomial
divisor'' clause of A2.EO.3c, and it is also the cleanest argument for why long division still earns
its keep after synthetic division shows up. \textbf{Part 2} inverts the Factor Theorem --- students
who understand it will set $f(3)=0$ immediately; students who do not will try to divide with an
unknown $k$ in the way. \textbf{Part 4 is the whole point of the tier:} groups discover for
themselves that the integer candidate list is incomplete, which is exactly the need the Rational
Root Theorem fills tomorrow. Do not resolve it for them --- let 4(d) be their conjecture.
\end{teachernote}

\end{document}
